% !TeX root = bifurcation-theory.tex
\section{bifurcation-theory}

\subsection[status=draft,tag=0000000,label=\labelNFM{subsec:family-of-dynamics-and-bifurcation-theorey}]{Family of dynamics and bifurcation theory}

Now let us turn to the family of dynamics.

Let $f: \bbP^d \times \Lambda \to \bbP^d \times \Lambda$ be a family of endomorphism on $\bbP^d$ over $\Lambda$.
Here $\Lambda$ may be an algebraic variety or a complex analytic variety.
We call that $f$ a holomorphic family or an algebraic family according to $\Lambda \to \End_{\bbP^d,q}(\bbC)$.

\begin{definition}[label=\labelNFM{def:relative-green-currents}, status=draft, tag=0000000]
	Let $f: \bbP^d \times \Lambda \to \bbP^d \times \Lambda$ be a holomorphic family of $q$-polarized endomorphism on $\bbP^d$ over $\Lambda$.
	Then the \emph{relative Green current} of $f$ is defined as
	\[ T_f \coloneqq \lim_{k \to +\infty} \frac{(f^k)^*\omega_2}{q^k} \]

	\Yang{}
\end{definition}

\begin{lemma}[label=\labelNFM{lem:relative-green-currents-and-green-currents}, status=draft, tag=0000000]
	Let $f: \bbP^d \times \Lambda \to \bbP^d \times \Lambda$ be a holomorphic family of $q$-polarized endomorphism on $\bbP^d$ over $\Lambda$.
	Then
	\[ T_f \wedge \delta_{\bbP^d \times \{\lambda\}} = T_{f_\lambda} \]
	\Yang{}
\end{lemma}

\begin{definition}[label=\labelNFM{def:bifurcation-currents}, status=draft, tag=0000000]
	Let $f: \bbP^d \times \Lambda \to \bbP^d \times \Lambda$ be a holomorphic family of $q$-polarized endomorphism on $\bbP^d$ over $\Lambda$.
	For every $\lambda \in \Lambda$, we define
	\[ L(\lambda) \coloneqq L(f_\lambda) \coloneqq \int_{\bbP^d} \log \|(\upd f_\lambda)_z^{\wedge d}\|_{\FS} \upd \mu_{f_\lambda}. \]
	The \emph{bifurcation currents} of $f$ is defined as
	\[ T_{\ bif} \coloneqq \ddc L \in \scrA^\Lambda_{-1,-1}(\Lambda) \]
	which is a currents on $\Lambda$ of degree $(1,1)$.
	\Yang{}
\end{definition}

\Yang{}

\begin{definition}[label=\labelNFM{def:bifurcation-currents-for-a-pair-af}, status=draft, tag=0000000]
	Let $f: \bbP^d \times \Lambda \to \bbP^d \times \Lambda$ be a holomorphic family of $q$-polarized endomorphism on $\bbP^d$ over $\Lambda$.
	For $a \in \bbP^d(\Lambda)$, we define the \emph{bifurcation measure} for the pair $(a,f)$ as
	\[ T_{f,a} \coloneqq (p_\Lambda)_*(T_f \wedge \delta_{\Gamma_a}) = a^*T_f. \]
	We say that $a$ is \emph{active} if $T_{f,a} \neq 0$.
	\Yang{}
\end{definition}


\begin{proposition}[label=\labelNFM{prop:bifurcation-currents-are-sum-of-critical-points}, status=draft, tag=0000000]
	Let $f: \bbP^1 \times \Lambda \to \bbP^1 \times \Lambda$ be a holomorphic family of $q$-polarized endomorphism on $\bbP^d$ over $\Lambda$.
	Then
	\[ T_{\ bif} = \sum_{c_i \in critical(f)} \mu_{f,c_i}. \]
	\Yang{}
\end{proposition}

%
%
% \begin{definition}[label=\labelNFM{def:bifurcation-measure}, status=draft, tag=0000000]
% 	\leavevmode\par
% 	\begin{latexmathchange}
% 		Let $f:\bbP^n\times\Lambda\to\bbP^n\times\Lambda$ be a holomorphic family of endomorphisms of algebraic degree $d\geq2$, where $\Lambda$ is a complex manifold of dimension $m\geq1$.
% 		Let $\mu_\lambda=T_{f_\lambda}^n$ be the equilibrium probability measure of $f_\lambda$.
% 		Write
% 		\[
% 			L(\lambda)\coloneqq\sum_{j=1}^n\chi_j(\lambda)
% 			=\int_{\bbP^n}\log\bigl\|\textstyle\bigwedge^n Df_\lambda(z)\bigr\|_{\mathrm{FS}}\,\upd\mu_\lambda(z),
% 		\]
% 		where $\chi_1(\lambda),\ldots,\chi_n(\lambda)$ are the complex Lyapunov exponents with respect to $\mu_\lambda$.
% 		The real-valued function $L$ is locally H\"older continuous and psh.
% 		The \emph{bifurcation current} of the family is
% 		\[
% 			T_{\mathrm{bif}}\coloneqq\upd_\Lambda\upd_\Lambda^c L.
% 		\]
% 		It is a closed positive current of bidegree $(1,1)$, or bidimension $(m-1,m-1)$, on $\Lambda$.
% 		Since $L$ is locally bounded, the Bedford--Taylor products $T_{\mathrm{bif}}^j$, $1\leq j\leq m$, are well-defined.
% 		We call these the \emph{higher bifurcation currents}, and use the normalization
% 		\[
% 			\mu_{\mathrm{bif}}\coloneqq T_{\mathrm{bif}}^m
% 		\]
% 		for the \emph{bifurcation measure}.
% 		See \cite[\S 2.5]{Dinh-Sibony-2010-DynamicsSeveralComplex} and, for rational maps, \cite{DeMarco-2003-DynamicsRationalMaps}.
% 	\end{latexmathchange}
% 	\begin{latexmathchange}
% 		For a holomorphic marked point $a:\Lambda\to\bbP^n$, set $a_j(\lambda)\coloneqq f_\lambda^j(a(\lambda))$.
% 		Its \emph{activity current} is the closed positive $(1,1)$-current
% 		\[
% 			T_{a,f}\coloneqq\lim_{j\to\infty}d^{-j}a_j^*\omega_{\mathrm{FS}},
% 		\]
% 		where the limit is weak and $\omega_{\mathrm{FS}}$ is the normalized Fubini--Study form on $\bbP^n$.
% 		When $m=1$, this is a measure, also denoted $\mu_{a,f}$.
% 		When $n=1$ and the critical points $c_1,\ldots,c_{2d-2}$ can be holomorphically marked, counted with multiplicity, we have
% 		\[
% 			T_{\mathrm{bif}}=\sum_{i=1}^{2d-2}T_{c_i,f}.
% 		\]
% 		See \cite[\S 4.2.3]{Berteloot-2013-BifurcationCurrentsHolomorphic}.
% 	\end{latexmathchange}
% 	% \toai{add definition of bifurcation currents} -- resolved.
% \end{definition}

\begin{theorem}[label=\labelNFM{thm:temporay}, status=draft, tag=0000000]
	Let $f$ be a holomorphic family of $q$-polarized dynamics and $a \in \bbP^1(\Lambda)$,
	TFAE:
	\begin{enumerate}
		\item $\mu_{a,f} = 0$;
		\item $a$ is stable.
	\end{enumerate}

	A point $a \in \bbP^1(\Lambda)$ is passive iff $a$ is preperiodic.
	\Yang{}
\end{theorem}






